% Build Part II v4.6 first, then this file, from formalization_v3/.
% Cross-document references use the companion Part II .aux file.
\documentclass[11pt,a4paper]{article}
\usepackage[T1]{fontenc}
\usepackage[utf8]{inputenc}
\usepackage{lmodern,microtype}
\usepackage[margin=.85in]{geometry}
\usepackage{amsmath,amssymb,amsthm,mathtools,aliascnt}
\usepackage{booktabs,tabularx,array,enumitem,needspace,float}
\usepackage[section]{placeins}
\usepackage{xcolor,tikz,mdframed,graphicx}
\usetikzlibrary{arrows.meta,positioning,calc,decorations.pathreplacing}
\usepackage{xr-hyper}
\usepackage[hidelinks]{hyperref}
\usepackage[nameinlink,noabbrev]{cleveref}
\usepackage{fancyhdr}
\hypersetup{pdftitle={Sync HotStuff RR, Part IV: Force Handler and Publish Layer (v0.1)},pdfsubject={Atomic branch updates, the Force runtime handler and Publish Layer implementation}}
\setlength{\headheight}{14pt}
\pagestyle{fancy}\fancyhf{}
\lhead{\small Sync HotStuff RR, Part IV}\rhead{\small Force and Publish Layer v0.1}
\cfoot{\thepage}
\setlength{\parindent}{0pt}\setlength{\parskip}{4pt}
\setlength{\emergencystretch}{2em}
\renewcommand{\topfraction}{.92}
\renewcommand{\textfraction}{.06}
\renewcommand{\floatpagefraction}{.85}
\makeatletter
\setlength{\@fptop}{0pt}\setlength{\@fpbot}{0pt plus 1fil}
\makeatother
\setlist{nosep,leftmargin=*}
\newcolumntype{Y}{>{\raggedright\arraybackslash}X}
\definecolor{rrRed}{HTML}{A51C30}
\definecolor{rrBlue}{HTML}{2563EB}
\definecolor{rrTeal}{HTML}{0F766E}
\newcommand{\RR}[1]{{\color{rrRed}#1}}
\newcommand{\QC}{\mathcal C}
\newcommand{\FC}{\mathcal F}
\newcommand{\Hash}{\mathsf H}
\newcommand{\rank}{\operatorname{rank}}
\newcommand{\TB}{\mathit{TB}}
\newcommand{\tagg}[1]{\ensuremath{\mathsf{#1}}}
\newcommand{\defterm}[1]{\textcolor{rrBlue}{\textit{#1}}}
\newtheorem{definition}{Definition}[section]
\newaliascnt{assumption}{definition}
\newtheorem{assumption}[assumption]{Assumption}
\aliascntresetthe{assumption}
\crefname{assumption}{assumption}{assumptions}
\Crefname{assumption}{Assumption}{Assumptions}
\newaliascnt{proposition}{definition}
\newtheorem{proposition}[proposition]{Proposition}
\aliascntresetthe{proposition}
\crefname{proposition}{proposition}{propositions}
\Crefname{proposition}{Proposition}{Propositions}
\newcounter{protocol}
% Protocol boxes float to the top of a page, so long boxes leave no gaps.
\newenvironment{protocol}[3][htbp]{\begin{figure}[#1]\refstepcounter{protocol}\label{#2}
\begin{mdframed}[linewidth=.6pt,skipabove=0pt,skipbelow=0pt,
innertopmargin=7pt,innerbottommargin=7pt]\small
\textbf{Protocol \theprotocol. #3}\par\smallskip}{\end{mdframed}\end{figure}}
\crefname{protocol}{Protocol}{Protocols}
\Crefname{protocol}{Protocol}{Protocols}
\title{Sync HotStuff RR, Part IV: Force Handler and Publish Layer (v0.1)}
\author{}\date{Version 0.1 --- 20 October 2026}
\externaldocument[II-]{sync_hotstuff_rr_part2_rotation_v4_6}
\begin{document}
\maketitle
\begin{abstract}
This document contains the Force handler and Publish Layer implementation
extracted from Part~II, version~4.6. It details how an explicit L1 atomic branch update closes
its target branch, installs a new region map and restarts from the
retained global prefix~$G$. The handler is the runtime form of the
atomic-update prelude in Part~II's global-finality rule; it does not
apply an update a second time. The Publish Layer implementation describes
leader-first submission, permissionless backup, fee sponsorship, retries
and retention under Part~II's delivery-service assumptions.
\end{abstract}

\paragraph{Companion document and scope.}
Read this document with
\href{sync_hotstuff_rr_part2_rotation_v4_6.pdf}{Part~II: Sequencer Rotation,
version~4.6}. Part~II defines the model, authorities, control log,
activation assumptions, replay state, target-index rule and Session
Manager. It retains the Force transition in the global-finality rule,
the handoff and Force-close interfaces, examples, the security and
lifecycle guarantees, and reference control policy. It also defines the
publication-service assumptions and source-retention requirements used
by the implementation here. All notation and assumptions here are those
of Part~II. References marked ``of Part~II'' point to that companion;
the remaining references are local to this document.

\section{The Force handler}
\label{sec:force}
\Cref{alg:force} details the atomic-update prelude of
\cref{II-alg:global-rule} of Part~II; it is the same transition, not a second
application of the update. It executes explicit L1 instructions.
Once the closure is verified, \cref{II-alg:rotation} of Part~II immediately terminates
the local old-version programs (\cref{II-sec:force-close} of Part~II). Silence, a
health report or a local timeout never runs this transition.

\subsection{What an instruction says}
\label{sec:force-stages}
An MSIG is an atomic branch update
$U=(\mathit{chain},c,\theta_a,\beta_{\rm target},S_{\rm new})$
(\cref{II-def:control-update} of Part~II). It names the branch version to close and the
complete map to install on the next version. These two effects are
inseparable: if the target is current, both occur; if the target is
stale, neither occurs. The update names no retained height. Where the new
branch starts is decided when the instruction executes: if the committee
saw $G=108$ when it signed but $G=112$ when the update takes effect, 112
is kept. The committee submits the instruction under
\cref{II-ass:committee,II-ass:services} of Part~II. L1 finalization makes it known; its
effect begins at $J_{q(\theta_a)}$. Until then, regular handoffs and
valid old-version output continue. A regular handoff changes $r$, not
$\beta$.

\paragraph{Why the target version is checked.}
Suppose instruction X says ``at time 120, close branch version 7''. If
in the meantime A only handed off to B by a regular handoff, the version
is still 7, so X executes at 120 and closes all of version 7, including
B's regular successor. But if another Force has already moved the
version to 8, X does not execute: its target is 7, and it never turns
into ``close whatever version is current''. So a late or duplicate
instruction cannot cut off a newer branch version: its target is stale
(\cref{II-def:control-effect} of Part~II).

\paragraph{The export session.}
The \defterm{export session} is responsible for the next DATA after $G$.
Its source context, fixed quota and accepted count decide the target
index of a Force (\cref{II-sec:retry-slot} of Part~II). It has no cross-region failure
timer. Its program may still be Assigned while this public responsibility
exists, and a later regular program may already be producing.

The newly installed map changes eligibility for the new branch only. An
update whose map leaves the source region ON can select that region again
when it is the only eligible one. If no region is eligible, the new
boundary keeps its anchor and target index without opening a regional
engine. Since the map is immutable within that version, only another
successful atomic update can supersede this unassigned boundary and try
a different complete map.

\begin{protocol}{alg:force}{Apply an atomic branch update}
Before the payload of authenticated finalized block $J_q$, let
$\theta=\theta_q$ and record the branch version $\beta_0$ at the start
of the block.
\begin{enumerate}[label=\textbf{\arabic*)},itemsep=3pt]
\item \textbf{Read the activated update.}
Use the authenticated activated-update input $\Lambda_q$ supplied for
$J_q$ (\cref{II-sec:logical-carrier} of Part~II). By \cref{II-ass:single-activation} of Part~II, it is empty
or a singleton. If it is empty, return to ordered payload processing.
Otherwise write
\[
 \Lambda_q=\{U_q\},\qquad
 U_q=(\mathit{chain},c_q,\theta_a,\beta_{\rm target},S_{\rm new}).
\]
Repeated source delivery never processes $U_q$ twice.

\item \textbf{Check its target.}
Advance the processed control cursor past $U_q$. If
$\beta_{\rm target}\ne\beta_0$, the whole update is stale: leave all
other protocol state unchanged and return to ordered payload processing. If
$\beta_{\rm target}=\beta_0$, continue to Step~3. No output, Ready,
elapsed-time or health condition is checked.

\item \textbf{Close, install the map and keep $G$.}
If $\mathsf{Tail}=\tagg{NEXT}(\nu)$, record
$\operatorname{sid}(\nu)$, $i(\nu)$ and $a(G,\nu)$; if
$\mathsf{Tail}=\tagg{PENDING}(\omega)$, read the stored source context
and target index from $\omega$.  Let $i_{\rm src}=i(\nu)$ in the first
case and the stored source region in the second. Keep all of $G$ as it
stands at the start of $J_q$, with its original block-bound records. Keep
$U_q$ and the header of $J_q$ as closure evidence. Close $\beta_0$, set
$\beta'=\beta_0+1$ exactly once and install
$S_{\beta'}:=S_{\rm new}$. Deliver this verified closure to the Manager's
Force-close rule immediately, independently of target-authority
availability. Fix the
target index $r'$ by \cref{II-sec:retry-slot} of Part~II, let
$A=\operatorname{last}(G)$ and compute
$i'=\tagg{NextLive}(i_{\rm src},S_{\beta'};\beta',r')$ using the new target's
$\tagg{AUTH}_{\beta',r',0}$ (\cref{II-sec:authority-context} of Part~II).
Reusing $r'$ does not reuse an old branch's authority. If required source
evidence is missing, pause without interpreting the candidate set. With complete evidence, set
\[
 \mathsf{Tail}:=
 \begin{cases}
 \tagg{PENDING}(\tagg{FORCE}(c_q,q),\beta',r',i_{\rm src},A,q),
   & i'=\bot,\\
 \tagg{NEXT}(\beta',r',i',A,h(A)+\TB_{r'},q,
              \tagg{FORCE}(c_q,q)),& i'\ne\bot.
 \end{cases}
\]
Archive the closed version's unconfirmed suffix and regular successors.
The new engine, when entered, starts with fresh core state and $g=t=0$.
No new-version DATA is accepted in $J_q$ itself.

\item \textbf{Persist once.}
Persist the processed control cursor, the new version and map, the closed
mark, $G$ and the new tail as one replay transition. Then return to
ordered payload processing.
\end{enumerate}
\end{protocol}

The cutoff and retained anchor are fixed even if target evidence arrives
later. Resume the same canonical operation; never rerun a persisted
update. A successful update can have a fixed closure while target
selection still waits for its authority evidence. In contrast, later
Espresso blocks alone never revisit a PENDING boundary: changing its map
requires another successful atomic update.

Every prelude executes at most one Force. The activation thresholds of
two Forces are at least $T_{\rm update}$ apart
(\cref{II-def:control-valid} of Part~II); no DATA, Ready or other output is required
between them.
Each update must still satisfy its own authorization,
advance-announcement and target-version conditions.

\subsection{Atomic maps, stale updates and the cutoff}
If an update's target is stale, the entire update is a no-op. It does not
install $S_{\rm new}$, close the current version, select a region or
choose an anchor. If the map change is still desired, the committee must
submit a fresh update targeting the current version. This atomic rule is
why a map change cannot retroactively revise a regular handoff within an
open version.

All control effects come before the payload of the triggering block.
DATA accepted in earlier blocks stays in $G$. An old-version result in
the triggering block cannot cancel a Force by arriving with late
progress. It may be kept as regional evidence but cannot append its
suffix. A newer version starts only from the fixed retained anchor.

For example, take threshold $d=120$ and Espresso blocks with timestamps
119, 122 and 125. The block at 119 accepts valid DATA of the current
version. The block at 122 is the first to reach 120, so the Force runs
first and old-version DATA in that block can no longer be accepted;
nor can any later old-version DATA. New-version DATA can be accepted only
from the block after 122.

A proof of closure contains the actual MSIG and the complete applicable
source context. A bare timestamp, an unavailable RPC response, an
unfinalized control or a monitoring alarm is not such a proof. After
learning of the closure, regional Session Managers apply
the Force-close rule of \cref{II-alg:rotation} of Part~II; the cutoff does not wait
for a node to learn it or release resources.

\section{Implementing the Publish Layer}
\label{sec:publish-implementation}
The protocols in Part~II use a delivery-service assumption. The following
leader-first design explains how existing participants can implement it.
It adds no consensus round and no exclusive publisher role.
The implementation must meet $T_{\rm pub}$; the sketch alone is not a
proof of that bound.

\subsection{Leader first, with permissionless backup}
The proposing regional leader is the preferred sender. Once it has the
complete finalized result, it builds and submits an EXPORT immediately.
Other replicas continue collecting FINAL messages, keeping payloads and
making complete results available to the service. Commitment and
Completion do not wait for a publication acknowledgement.

Another replica or an external data verifier may submit the same result
without the leader's permission. An implementation may give the leader a
short first opportunity, but the backup delay, data retrieval and retries
must fit within $T_{\rm pub}$ from complete-result availability. A private
announcement that the leader has sent something is not enough to suppress
backup indefinitely. The service checks submission and Espresso evidence.

The fallback sender uses the original DATA and receipt. It cannot replace
the DATA, rewrite its session pair or create FINAL on behalf of the region.
The service guarantees delivery through an available path; no particular
sender is guaranteed to win the submission race.

\subsection{Fee sponsorship}
A backup sender need not hold tokens or pay from its own balance. It may
use a funded relay, a prepaid submission account, or another sponsored
route that can submit to Espresso. The route carries the verified result;
it does not acquire authority to change its content.

Someone still supplies the fee resources. To justify the publication
bound, a complete result covered by the service assumption must have access to a funded,
working route even when the original leader is unavailable. A sender can
try another route if one refuses or fails. The concrete sponsorship and
admission mechanism is an implementation choice; this protocol assumes
no Espresso application smart contract for it.

\subsection{Duplicates, retries and acceptance}
Different senders may carry the same complete result in different EXPORTs.
The signed producing pair and block hash identify each DATA; changing a
receipt encoding, signer set or outer hint does not create a new output.
Accepted ancestors are reused as evidence and receive no second count.
A compatible next DATA can be extracted even when the sender's expected
base is stale or its result is listed only among supporting ancestry.

The Session Manager tracks the original EXPORT record and its per-DATA
acceptance. A complete next result already published with a bad hint is
interpreted at that source position under the same rule; a corrected
EXPORT is unnecessary.
If a genuine gap or an absent required object prevents acceptance, the
objects already published stay in $\mathcal O_q$; once the missing
objects appear in a later block, \cref{II-alg:global-rule} of Part~II accepts the
result there without republication. No bound on when the missing
objects arrive is inferred. A closed version never regains append
permission.

The two service bounds concern complete-result delivery and visibility.
The conditional progress discussion of \cref{II-sec:global-progress} of Part~II
uses those actual recorded bytes. It does not assume a bounded retry
with a corrected expected-base hint. Terminal DATA commits before successor authority
resolution. Preparation for a selected live region has the separate
scope of \cref{II-ass:ready-service} of Part~II; arbitrary archive catch-up has no
uniform bound. \Cref{II-sec:transport-scope} of Part~II states what the transport
must put into Espresso for the rule to use it.

\subsection{Retention and incentives}
Before finalized Espresso inclusion, the holders of a complete result
keep its payload and evidence available to backup senders. Afterward,
Espresso's retained, authenticated history supplies returning readers and
initial hydration. Prepared regional nodes additionally keep the raw
source bytes under \cref{II-sec:source-retention} of Part~II; finalization does not
replace that retained copy by an eventual archive pointer. Direct transfer
can still accelerate regular handoff and state recovery. Program
termination and checkpointing preserve required evidence and outgoing
sends.

Submission rewards and reimbursement may fund this service. They must not
count duplicate DATA as new progress. A report of delayed delivery
alone does not prove that the original leader, verifier or sponsor
equivocated. Health policy and any objectively verifiable penalties
are separate from the Force acceptance rule.

\end{document}
